%! TEX root = ../am1-20-21.tex
\chapter{Șiruri de numere reale și pozitive}

\section{Exerciții recapitulative}

În toate cele de mai jos, presupunem că lucrăm cu șiruri de numere reale și pozitive.

\vspace{1cm}

1. Găsiți valoarea de adevăr a afirmațiilor de mai jos. Dacă sînt adevărate,
justificați. Dacă sînt false, găsiți un contraexemplu.
\begin{enumerate}[(a)]
    \item Orice șir monoton este mărginit.
    \item Orice șir mărginit este monoton.
    \item Orice șir convergent este monoton.
    \item Orice subșir al unui șir monoton este monoton.
    \item Suma a două șiruri monotone este un șir monoton.
    \item Orice șir divergent este nemărginit.
    \item Dacă șirul format din pătratele termenilor unui șir este convergent,
        atunci șirul inițial este convergent.
\end{enumerate}

\vspace{1cm}

2. Calculați limitele șirurilor cu termenul general $ a_n $ în cazurile de mai jos:
\begin{enumerate}[(a)]
    \item $ a_n = \dfrac{n^3 + 5n^2 + 1}{6n^3 + n + 4} $;
    \item $ a_n = \dfrac{n^2 + 2n + 5}{5n^3 - 1} $;
    \item $ a_n = \arcsin \dfrac{n + 1}{2n + 3} $;
    \item $ a_n = \dfrac{\sqrt{n^2 + 1} + 4\sqrt{n^2 + n + 1}}{n} $;
    \item $ a_n = \dfrac{1^2 + 2^2 + 3^2 + 4^2 + \dots + n^2}{n^3} $;
    \item $ a_n = \dfrac{1}{1 \cdot 3} + \dfrac{1}{3 \cdot 5} + \dots + \dfrac{1}{(2n - 1)(2n + 1)} $;
    \item $ a_n = \sqrt[3]{n^2 + 1} - \sqrt[3]{n^2 + n} $;
    \item $ a_n = 1 + \dfrac{1}{2} + \dfrac{1}{3} + \dots + \dfrac{1}{n} $;
    \item $ a_n = \left( \dfrac{n+5}{3n + 2} \right)^n $;
    \item $ a_n = 1 + \dfrac{1}{4} + \dfrac{1}{16} + \dfrac{1}{64} + \dots + \dfrac{1}{4^n} $;
    \item $ a_n = \dfrac{2 + \sin n}{n^2} $.
\end{enumerate}
